\documentclass[../../../../main.tex]{subfiles}
\begin{document}

\begin{flushright}
\footnotesize\textit{【自動文字起こし・要確認】}
\end{flushright}

[解] 到達時刻 $t_0$ とする。右のように点 P, Q を定める。
\[ |\overline{OP}| = |t-a|, \quad |\overline{OQ}| = |t_0-b|, \quad |\overline{PQ}| = \sqrt{2}(t_0-t) \]
だから、$\triangle OPQ$ にピタゴラスの定理を用いて、
\[ 2(t_0-t)^2 = (t-a)^2 + (t_0-b)^2 \quad \cdots \text{①} \]

\begin{tikzpicture}
\draw[->] (-1,0) -- (4,0) node[right] {$x$};
\draw[->] (0,-1) -- (0,4) node[above] {$y$};
\coordinate (O) at (0,0);
\coordinate (P) at (2,0);
\coordinate (Q) at (0,3);
\draw (P) -- (Q);
\node[below left] at (O) {O};
\node[below right] at (P) {P};
\node[above left] at (Q) {Q};
\node[below] at (2,-0.2) {$t-a$};
\node[left] at (-0.2,3) {$t_0-b$};
\node[above right] at (1,1.5) {$\sqrt{2}(t_0-t)$};
\end{tikzpicture}

$S = t_0 - t$ とすると、$S$ を min にする $t$ をもとめれば良い。①に代入して
\begin{align*}
2 S^2 &= (t-a)^2 + (S + t - b)^2 \\
&= S^2 + 2(t-b)S + (t-b)^2 + (t-a)^2
\end{align*}
\[ \therefore S^2 - 2(t-b)S - (t-a)^2 - (t-b)^2 = 0 \]
\[ \therefore S = t-b + \sqrt{2(t-b)^2 + (t-a)^2} \quad (\because S > 0) \]
であるから、両辺 $t$ で微分して、
\[ \frac{dS}{dt} = 1 + \frac{2(t-b) + (t-a)}{\sqrt{2(t-b)^2 + (t-a)^2}} \]
だから、
\[ \frac{dS}{dt} \ge 0 \iff \sqrt{2(t-b)^2 + (t-a)^2} \ge -2(t-b) - (t-a) \]
$-2(t-b) - (t-a) \ge 0 \iff t \le \frac{a+2b}{3}$ のもとで2乗して解く（$t \ge \frac{a+2b}{3}$ のときは常に $\frac{dS}{dt} \ge 0$ ）
\begin{align*}
2(t-b)^2 + (t-a)^2 &\ge 4(t-b)^2 + (t-a)^2 + 4(t-a)(t-b) \\
\iff 0 &\ge 2(t-b)^2 + 4(t-b)(t-a) \\
&= 2(t-b) \{ 3t - (2a+b) \} \\
\iff \frac{2a+b}{3} &\le t \le b \quad (\because 0 < a < b)
\end{align*}
だから、下表をうる。

\begin{center}
\begin{tabular}{|c|c|c|c|c|c|}
\hline
$t$ & $0$ & $\cdots$ & $\frac{2a+b}{3}$ & $\cdots$ & $b$ \\
\hline
$S'$ & & $-$ & $0$ & $+$ & \\
\hline
$S$ & & $\searrow$ & & $\nearrow$ & \\
\hline
\end{tabular}
\end{center}
したがって、もとめるのは $t = \frac{2a+b}{3}$ である。

\end{document}
